% ==========================================
% LatexRefinement Step Output: step4-05-26-25-tuesday-all-offset-final.tex
% Model: gemini-3.5-flash
% Temperature: 0.33
% TopP: 0.95
% TopK: 40
% MaxOutputTokens: 65535
% ThinkingBudget: 32765
% ThinkingLevel: HIGH
% Prompt Tokens: 45'305
% Candidates Tokens: 19'083 (inkl. Thinking Tokens)
% Cached Tokens: 0
% Processed on: 2026-07-18 14:21:46
% ==========================================

% Primary Language: German

\lecturechapter[14]{Dienstag}{26. Mai}{26. Mai 2026}{Polynomringe, Restklassenkörper und endliche Körper}
% Old chapter title: Polynomringe und Division mit Rest

\begin{nice-box}[Vorlesungsübersicht \& Kontext]
In dieser letzten Vorlesung der Reihe vertiefen wir unser Verständnis der Algebra anhand von Polynomringen über Körpern. Wir untersuchen die strukturellen Analogien zum Ring der ganzen Zahlen, insbesondere die Division mit Rest, Teilbarkeit und irreduzible Polynome. Darauf aufbauend konstruieren wir Restklassenkörper und klassifizieren endliche Körper. Zum Abschluss besprechen wir organisatorische Details und Vorbereitungstipps für die bevorstehende Prüfung.

\textbf{Themenschwerpunkte:}
\begin{itemize}
    \item \textbf{Polynomringe:} Definition, Gradfunktion und fundamentale Eigenschaften über Körpern.
    \item \textbf{Division mit Rest \& Teilbarkeit:} Euklidischer Algorithmus, größter gemeinsamer Teiler ($\operatorname{ggT}$) und irreduzible Polynome.
    \item \textbf{Restklassenringe:} Konstruktion von Körpern (wie $\mathbb{F}_4$) und der Klassifikationssatz endlicher Körper.
    \item \textbf{Prüfungsvorbereitung:} Format, Hilfsmittel und wertvolle Tipps für eine erfolgreiche Prüfung.
\end{itemize}
\end{nice-box}

\section{Einführung und Polynomringe}

\begin{spoken-clean}[00:00:00 - 00:00:38]
Hallo zusammen und herzlich willkommen. Wir fangen an. 

Schön, dass Sie so zahlreich erschienen sind, trotz der Hitze hier Algebra zu lernen, anstatt sich im Pool zu erfrischen. Wir werden heute, wie gesagt, noch Algebra machen --- wir schauen uns noch ein bisschen Polynomringe an und insbesondere noch endliche Körper. Auch das wird wieder nur so ein bisschen ein Vorgeschmack auf Algebra. Und dann haben wir hoffentlich am Ende noch etwas Zeit, um kurz über die Prüfung zu sprechen. Gut.
\end{spoken-clean}

\begin{spoken-clean}[00:00:38 - 00:02:27]
Aber jetzt noch zurück. Wir hatten letzte Woche etwas elementare Zahlentheorie gesehen und haben mit dem Ring der ganzen Zahlen gearbeitet. Jetzt schauen wir uns noch etwas mehr im Detail einen weiteren wichtigen Ring an: das ist der Ring der Polynome über einem Körper. Auch hier gibt es einfach nur einen kleinen Einblick, was es da so gibt. 

Also, wenn $K$ ein Körper ist, dann können wir die Menge der Polynome über $K$ anschauen --- also alle formalen Summen von der Form $a_i x^i$ (d.\,h. $a_i$ mal $x$ hoch $i$) für $i$ gleich $0$ bis $n$ für ein $n \ge 0$, wobei die $a_i$ alle Elemente in $K$ sind. Das ist die Menge der Polynome über $K$ mit der üblichen Multiplikation und Addition.
\end{spoken-clean}

\subsection{Definition des Polynomrings}

\begin{math-stroke}
\begin{definition}[Polynomring]\label[definition]{def:polynomring}
Sei $K$ ein Körper. Der \newterm{Polynomring} in einer Variablen $x$ über $K$, bezeichnet als $K[x]$, ist definiert als die Menge aller formalen Summen:
\[
K[x] := \left\{ \sum_{i=0}^n a_i x^i \;\middle|\; n \in \mathbb{N}, a_i \in K \right\}
\]
versehen mit der üblichen Addition und Multiplikation von Polynomen.
\end{definition}
\end{math-stroke}

\begin{spoken-clean}[00:02:27 - 00:06:20]
Und das ist der Polynomring in einer Variablen über $K$. Gut, ich gehe davon aus, dass Sie das bereits in der Linearen Algebra oder Analysis gesehen haben. Das ist wichtig: Polynomringe sind sehr wichtige Beispiele von Ringen. Und das Wichtige, das man wissen muss, ist: Es sind wirklich einfach formale Summen von dieser Form. Also ist der Polynomring nicht zu verwechseln mit dem Ring der polynomialen Funktionen.

Das haben Sie hoffentlich auch gesehen; aber einfach, um das nochmals klarzumachen: Also ist $K[x]$ im Allgemeinen nicht dasselbe wie der Ring der polynomialen Abbildungen von $K$ nach $K$. Eine polynomiale Abbildung ist eine Abbildung, die durch ein Polynom gegeben ist --- das heißt, Abbildungen der Form, bei denen $x$ nach $p(x)$ geschickt wird, für ein Polynom $p(x)$ in $K[x]$.
\end{spoken-clean}

\begin{math-stroke}[Polynomring vs. Polynomiale Abbildungen]
\begin{remark}[Polynomiale Abbildungen]\label[remark]{rem:polynomial-mappings}
Der Polynomring $K[x]$ ist im Allgemeinen nicht identisch mit dem Ring der polynomialen Abbildungen von $K$ nach $K$, d.\,h. Abbildungen der Form:
\[
K \to K, \quad x \mapsto p(x) \quad \text{für ein} \quad p(x) \in K[x].
\]
\end{remark}

\begin{example}[Endliche Körper]\label[example]{ex:finite-fields}
Sei $K = \mathbb{F}_p$ mit einer Primzahl $p$. Dann ist der Polynomring $K[x]$ unendlich (da es Polynome beliebig hohen Grades gibt), während die Menge aller Abbildungen von $K$ nach $K$ endlich ist (es gibt genau $p^p$ solche Abbildungen). Insbesondere induzieren unendlich viele verschiedene Polynome dieselbe Abbildung.
\end{example}
\end{math-stroke}

\begin{spoken-clean}[00:06:20 - 00:08:05]
Also jedes Polynom induziert eine polynomiale Abbildung von $K$ nach $K$, halt einfach gegeben durch $x \mapsto p(x)$. Aber es gibt natürlich a priori --- insbesondere wenn $K$ endlich ist --- viele verschiedene Polynome, die dieselbe Abbildung induzieren. 

Denn wie wir im Beispiel gesehen haben, gibt es endliche Körper. Wenn jetzt unser $K$ zum Beispiel der Körper $\mathbb{F}_p$ ist --- mit einer Primzahl $p$ ---, dann ist der Ring der Polynome über $K$ natürlich unendlich, weil man hier einen beliebig hohen Grad wählen kann und es somit unendlich viele Polynome gibt. Aber es gibt nur endlich viele Abbildungen von $K$ nach $K$.
\end{spoken-clean}

\begin{spoken-clean}[00:08:05 - 00:10:45]
Das heißt insbesondere, es gibt nur endlich viele polynomiale Abbildungen. Das ist einfach noch eine wichtige Bemerkung, damit man das nicht verwechselt. 

Ich glaube, in der Mittelschule ist man da oft ungenau, weil man meistens über dem Körper der reellen Zahlen arbeitet --- da gibt es kein Problem, da kann man das direkt so identifizieren. Aber sobald man über einem endlichen Körper arbeitet, ist der Polynomring natürlich viel größer als der Ring der polynomialen Abbildungen.
\end{spoken-clean}

\subsection{Der Grad eines Polynoms}

\begin{spoken-clean}[00:10:45 - 00:12:44]
Machen wir noch eine Definition, die Sie auch schon gesehen haben. Wenn wir ein Polynom $f$ der Form $a_0 + a_1 x + \dots + a_n x^n$ haben, sodass $a_n \neq 0$ ist und danach alle weiteren Koeffizienten null sind, dann ist der Grad von $f$ genau dieses $n$ --- also das größte $n$, sodass der Koeffizient vor $x^n$ ungleich null ist. 

Das ist natürlich nur definiert, wenn $f$ nicht das Nullpolynom ist. Denn wenn $f$ das Nullpolynom ist, dann gibt es keinen Koeffizienten ungleich null --- alle Koeffizienten sind null. Das heißt, in diesem Fall müssen wir separat definieren, was der Grad ist: Wir setzen den Grad des Nullpolynoms als minus unendlich ($-\infty$) fest. Und somit ist der Grad für alle Polynome definiert. Ja, das wirkt im ersten Moment vielleicht etwas unglücklich, ist aber durchaus sinnvoll.
\end{spoken-clean}

\begin{math-stroke}
\begin{definition}[Grad eines Polynoms]\label[definition]{def:grad-polynom}
Sei $f = a_0 + a_1 x + \dots + a_n x^n \in K[x]$ mit $a_n \neq 0$. Der \newterm{Grad} von $f$, bezeichnet als $\deg(f)$, ist definiert als:
\[
\deg(f) := n.
\]
Für das Nullpolynom definieren wir:
\[
\deg(0) := -\infty.
\]
\end{definition}
\end{math-stroke}

\begin{spoken-clean}[00:12:44 - 00:13:35]
Sie sehen beispielsweise schon im nächsten Lemma, dass diese Definition sinnvoll ist. Es gilt Folgendes: Wenn wir zwei Polynome $f$ und $g$ in einer Variablen haben, dann ist der Grad von $f \cdot g$ gleich dem Grad von $f$ plus dem Grad von $g$. Das heißt, der Grad verhält sich additiv unter Multiplikation.
\end{spoken-clean}

\begin{math-stroke}[Grad eines Produkts]
\begin{lemma}\label[lemma]{lem:grad-produkt}
Für beliebige Polynome $f, g \in K[x]$ gilt:
\[
\deg(f \cdot g) = \deg(f) + \deg(g).
\]
\end{lemma}

\begin{short-proof}[Beweis der Gradaussagen]
Falls $f = 0$ oder $g = 0$ gilt, ist die Aussage mit der Konvention $-\infty + d = -\infty$ für alle $d \in \mathbb{N} \cup \{-\infty\}$ trivialerweise erfüllt.

Seien nun $f, g \neq 0$ mit:
\begin{align*}
f &= a_0 + a_1 x + \dots + a_n x^n \quad \text{mit} \quad a_n \neq 0 \implies \deg(f) = n, \\
g &= b_0 + b_1 x + \dots + b_m x^m \quad \text{mit} \quad b_m \neq 0 \implies \deg(g) = m.
\end{align*}
Das Produkt $f \cdot g$ lässt sich schreiben als:
\[
f \cdot g = a_0 b_0 + (a_0 b_1 + a_1 b_0) x + \dots + a_n b_m x^{n+m}.
\]
Da $K$ ein Körper ist (und somit insbesondere nullteilerfrei), folgt aus $a_n \neq 0$ und $b_m \neq 0$ direkt:
\[
a_n b_m \neq 0.
\]
Der Koeffizient des Terms mit der höchsten Potenz $x^{n+m}$ ist somit ungleich Null, woraus folgt:
\[
\deg(f \cdot g) = n + m = \deg(f) + \deg(g).
\]
\end{short-proof}
\end{math-stroke}

\begin{spoken-clean}[00:13:35 - 00:15:00]
Und zum Beweis: Zuerst sehen wir hier, dass unsere Definition Sinn ergibt, falls $f$ oder $g$ gleich $0$ ist. Es passt zusammen, wenn wir vereinbaren, dass $-\infty$ plus irgendeine Zahl immer noch $-\infty$ ist. Deswegen ist es sinnvoll, den Grad von $0$ als $-\infty$ zu definieren, damit die Formel auch in diesem Fall stimmt. 

Ansonsten sieht man das direkt. Ich schreibe das jetzt nicht allzu detailliert auf: Wenn wir $f = a_0 + a_1 x + \dots + a_n x^n$ und $g = b_0 + b_1 x + \dots + b_m x^m$ schreiben, wobei $a_n \neq 0$ und $b_m \neq 0$ sind, dann ist der Grad von $f$ gleich $n$ und der Grad von $g$ gleich $m$. 

Wenn man das Produkt $f \cdot g$ ausmultipliziert, erhält man als konstanten Term $a_0 b_0$. Dann muss man beim Aufaddieren der Koeffizienten etwas aufpassen --- der lineare Term ist $(a_0 b_1 + a_1 b_0)x$ ---, und wenn man so weitergeht, ist der Term von höchster Ordnung schließlich $a_n b_m x^{n+m}$. 

Wichtig ist jetzt: Da $K$ ein Körper ist, ist dieser Leitkoeffizient ungleich null. Da $a_n \neq 0$ und $b_m \neq 0$ sind, ist auch ihr Produkt $a_n b_m \neq 0$. Somit folgt, dass der Grad von $f \cdot g$ gleich $n + m$ ist, was genau der Summe der Grade von $f$ und $g$ entspricht. Das ist sehr direkt. Aber eben...
\end{spoken-clean}

\begin{spoken-clean}[00:15:00 - 00:18:00]
Aufpassen muss man hier wirklich: Wir müssen voraussetzen, dass wir über einem Körper arbeiten. Man kann Polynome natürlich auch über beliebigen Ringen definieren, das haben Sie bestimmt auch schon getan. Aber wenn Sie sich das zum Beispiel über dem Restklassenring $\mathbb{Z}/4\mathbb{Z}$ anschauen, dann wissen Sie: $[2]$ ist ungleich $[0]$, aber $[2] \cdot [2]$ ist gleich $[0]$. Da haben Sie also Nullteiler. 

Und dann kann es durchaus vorkommen, dass der Grad von $f \cdot g$ strikt kleiner ist als die Summe der Grade von $f$ und $g$. Über beliebigen Ringen gilt im Allgemeinen nur, dass der Grad von $f \cdot g$ kleiner gleich dem Grad von $f$ plus dem Grad von $g$ ist. Aber bei uns gibt es keine Nullteiler, alles ist wunderbar --- das heißt, wir haben tatsächlich Gleichheit.
\end{spoken-clean}

\begin{didactic-insight}[Gradformel über Ringen mit Nullteilern]
Die Additivität des Grades gilt im Allgemeinen nur über Integritätsbereichen (und damit insbesondere über Körpern). Besitzt ein Ring $R$ Nullteiler $a, b \neq 0$ mit $a \cdot b = 0$, so kann das Produkt der Leitterme zweier Polynome verschwinden, wodurch sich der Grad des Produktpolynoms verringert. In diesem Fall gilt lediglich die Abschätzung:
\[
\deg(f \cdot g) \le \deg(f) + \deg(g).
\]
\end{didactic-insight}

\section{Division mit Rest}

\begin{spoken-clean}[00:18:00 - 00:22:50]
Solche Fragen werden Sie in Ihren Algebra-Vorlesungen vorwärts und rückwärts durchstudieren; aber hier, erst einmal über Körpern, ist alles wunderbar. Gut. 

Das Schöne ist: Ein Polynomring in einer Variablen über einem Körper verhält sich in vielerlei Hinsicht sehr ähnlich wie der Ring der ganzen Zahlen. Das ist ein bisschen die Philosophie der heutigen Stunde. Zuerst einmal werden wir Division mit Rest machen. Das kennen Sie wahrscheinlich aus der Mittelschule, oder? Da haben Sie wahrscheinlich Polynomdivision gemacht und geschaut, was für ein Rest übrig bleibt. Und das kann man tatsächlich über beliebigen Körpern machen. 

Wenn $f$ und $g$ Polynome in einer Variablen sind (mit $f \neq 0$), dann gibt es eindeutige Polynome $q$ und $r$, sodass wir schreiben können: $g = q \cdot f + r$. Wir können also $g$ durch $f$ mit Rest teilen. Bei den ganzen Zahlen war das ja so, dass der Rest im Betrag kleiner als der Divisor war --- beziehungsweise wir haben gefordert, dass er positiv und kleiner als der Divisor ist. 

Das Schöne ist: Im Polynomring haben wir ebenfalls eine passende Größenfunktion, nämlich den Grad. Da können wir einfach den Grad nehmen. Das heißt, für den Rest $r$ fordern wir, dass der Grad von $r$ strikt kleiner als der Grad von $f$ ist. Okay, und das ist die Division mit Rest. Wir können uns den Beweis kurz anschauen.
\end{spoken-clean}

\begin{math-stroke}[Satz über die Division mit Rest]
\begin{theorem}[Division mit Rest]\label[theorem]{thm:division-mit-rest}
Seien $f, g \in K[x]$ mit $f \neq 0$. Dann existieren eindeutige Polynome $q, r \in K[x]$, sodass:
\[
g = q \cdot f + r \quad \text{mit} \quad \deg(r) < \deg(f).
\]
\end{theorem}
\end{math-stroke}

\begin{spoken-clean}[00:22:50 - 00:24:30]
Zuerst beweisen wir die Existenz, danach die Eindeutigkeit. Den Existenzbeweis führen wir genau so, wie man das aus der Schule kennt: per Induktion über den Grad von $g$. 

Falls der Grad von $g$ kleiner als der Grad von $f$ ist, ist es klar: Dann können wir einfach $q = 0$ und $r = g$ wählen und schreiben $g = 0 \cdot f + g$. Das ist kein Problem. 

Falls der Grad von $g$ größer gleich dem Grad von $f$ ist, machen wir den Induktionsschritt. Wir nehmen an, die Aussage sei bereits für alle Polynome mit kleinerem Grad gezeigt, und betrachten nun ein $g$ mit Grad $n \ge \deg(f)$.
\end{spoken-clean}

\begin{spoken-clean}[00:24:30 - 00:28:00]
Wir schreiben $g = a_0 + a_1 x + \dots + a_n x^n$ und $f = b_0 + b_1 x + \dots + b_m x^m$. Wir wissen, dass $n \ge m$ ist, und die Leitkoeffizienten $a_n$ und $b_m$ sind ungleich null. Da wir über einem Körper arbeiten, existiert insbesondere das Inverse $b_m^{-1}$. 

Wir definieren nun ein neues Polynom $g_1$ als $g$ minus einen Korrekturterm, um den Leitterm von $g$ zu eliminieren. Das heißt, wir subtrahieren $b_m^{-1} a_n x^{n-m} f$. 

Wenn wir das ausschreiben, sehen wir, dass der highest Term von $b_m^{-1} a_n x^{n-m} f$ genau $a_n x^n$ ist. Dieser hebt den Leitterm $a_n x^n$ von $g$ exakt auf. In der Differenz kürzt sich der Term mit $x^n$ also weg, und was übrig bleibt, ist ein Polynom $g_1$ mit einem Grad, der strikt kleiner als $n$ ist. Der Grad von $g_1$ ist also strikt kleiner als der Grad von $g$.
\end{spoken-clean}

\begin{proof}[Beweis der Division mit Rest]
\begin{math-stroke}[Existenzbeweis]
Wir beweisen die Existenz der Darstellung mittels starker Induktion über den Grad $n := \deg(g) \in \mathbb{N} \cup \{-\infty\}$:

\textbf{Induktionsanfang:}
Falls $\deg(g) < \deg(f)$ gilt, wählen wir:
\[
q = 0 \quad \text{und} \quad r = g.
\]
Dann ist $g = 0 \cdot f + g$ mit $\deg(r) = \deg(g) < \deg(f)$ eine gültige Darstellung.

\textbf{Induktionsschritt:}
Es gelte nun $\deg(g) \ge \deg(f)$. Wir nehmen an, die Existenz einer solchen Darstellung sei bereits für alle Polynome mit einem Grad strikt kleiner als $n = \deg(g)$ bewiesen.

Seien die Polynome gegeben durch:
\begin{align*}
g &= a_0 + a_1 x + \dots + a_n x^n \quad \text{mit} \quad a_n \neq 0 \implies \deg(g) = n, \\
f &= b_0 + b_1 x + \dots + b_m x^m \quad \text{mit} \quad b_m \neq 0 \implies \deg(f) = m.
\end{align*}
Da $\deg(g) \ge \deg(f)$ vorausgesetzt ist, gilt $n \ge m$. Da $K$ ein Körper ist, existiert das multiplikative Inverse $b_m^{-1} \in K$. Wir definieren ein neues Polynom $g_1 \in K[x]$ durch:
\begin{equation}\label{eq:g1-definition}
g_1 := g - b_m^{-1} a_n x^{n-m} f.
\end{equation}
Wir betrachten den Koeffizienten von $x^n$ in $g_1$:
\[
g_1 = \left( a_0 + \dots + a_n x^n \right) - b_m^{-1} a_n x^{n-m} \left( b_0 + \dots + b_m x^m \right).
\]
Der Term von höchster Ordnung im subtrahierten Teil ist genau:
\[
b_m^{-1} a_n x^{n-m} \cdot b_m x^m = a_n x^n.
\]
Dieser hebt den Leitterm $a_n x^n$ von $g$ exakt auf. Daraus folgt unmittelbar:
\[
\deg(g_1) < \deg(g) = n.
\]
\end{math-stroke}

\begin{spoken-clean}[00:28:00 - 00:28:47]
Da der Grad von $g_1$ kleiner ist, können wir die Induktionshypothese anwenden. Es existieren also Polynome $q_1, r_1 \in K[x]$, sodass wir $g_1 = q_1 \cdot f + r_1$ schreiben können, wobei der Grad von $r_1$ strikt kleiner als der Grad von $f$ ist. 

Das heißt, wir können $g$ schreiben als $g = (b_m^{-1} a_n x^{n-m} + q_1) f + r_1$. Wenn wir nun $q := b_m^{-1} a_n x^{n-m} + q_1$ und $r := r_1$ definieren, erhalten wir die gewünschte Darstellung $g = q \cdot f + r$. 

Damit ist die Existenz gezeigt. Die Eindeutigkeit beweisen wir hier nicht im Detail, da der Beweis vollkommen analog zum Fall der ganzen Zahlen $\mathbb{Z}$ verläuft.
\end{spoken-clean}

\begin{math-stroke}[Abschluss des Existenzbeweises]
Da $\deg(g_1) < \deg(g)$, können wir die Induktionsvoraussetzung auf $g_1$ anwenden. Es existieren somit Polynome $q_1, r_1 \in K[x]$ mit:
\[
g_1 = q_1 \cdot f + r_1 \quad \text{mit} \quad \deg(r_1) < \deg(f).
\]
Setzen wir dies in Gleichung \eqref{eq:g1-definition} ein, erhalten wir:
\begin{align*}
g - b_m^{-1} a_n x^{n-m} f &= q_1 \cdot f + r_1 \\
\implies g &= \left( b_m^{-1} a_n x^{n-m} + q_1 \right) f + r_1.
\end{align*}
Mit den Definitionen:
\[
q := b_m^{-1} a_n x^{n-m} + q_1 \quad \text{und} \quad r := r_1
\]
ergibt sich die gewünschte Darstellung $g = q \cdot f + r$ mit $\deg(r) < \deg(f).$

\textbf{Eindeutigkeit:}
Der Beweis der Eindeutigkeit verläuft vollkommen analog zum Fall der ganzen Zahlen $\mathbb{Z}$ unter Ausnutzung der Gradformel (\cref{lem:grad-produkt}).
\end{math-stroke}
\end{proof}

\section{Teilbarkeit und der grösste gemeinsame Teiler}

\begin{math-stroke}[Teilbarkeit in K[x]]
\begin{definition}[Teilbarkeit]\label[definition]{def:teilbarkeit-polynom}
Seien $f, g \in K[x]$. Wir sagen $g$ \newterm{teilt} $f$ (oder $g$ ist ein \newterm{Teiler} von $f$), falls ein Polynom $h \in K[x]$ existiert, sodass:
\[
f = g \cdot h.
\]
Wir schreiben dafür $g \mid f$.
\end{definition}
\end{math-stroke}

\begin{math-stroke}[Der grösste gemeinsame Teiler]
\begin{theorem}[Existenz des ggT und Bezout-Identität]\label[theorem]{thm:ggt-bezout}
Seien $f, g \in K[x]$. Dann existiert ein Polynom $h \in K[x]$ mit den Eigenschaften:
\begin{enumerate}[label=\textbf{(\alph*)}]
    \item $h \mid f$ und $h \mid g$ (gemeinsamer Teiler),
    \item für jedes Polynom $p \in K[x]$ mit $p \mid f$ und $p \mid g$ gilt $p \mid h$.
\end{enumerate}
Zudem existieren Polynome $A, B \in K[x]$, sodass gilt:
\[
A \cdot f + B \cdot g = h.
\]
Ein solches Polynom $h$ heißt \newterm{grösster gemeinsamer Teiler} von $f$ und $g$.
\end{theorem}

\begin{remark}[Eindeutigkeit des ggT]
Der grösste gemeinsame Teiler zweier Polynome ist im Allgemeinen nur eindeutig bis auf Multiplikation mit Einheiten (d.\,h. konstanten Polynomen ungleich Null in $K^*$). Er ist eindeutig bestimmt, wenn man zusätzlich fordert, dass er normiert ist (Leitkoeffizient gleich $1$).
\end{remark}
\end{math-stroke}

\subsection{Eindeutigkeit des \texorpdfstring{$\operatorname{ggT}$}{ggT}}

\begin{spoken-clean}[00:28:47 - 00:30:06]
Man muss hier allerdings aufpassen: Bei den ganzen Zahlen ist der $\operatorname{ggT}$ eindeutig, wenn wir fordern, dass er positiv sein soll. Dort hat man den Vorteil einer kanonischen Wahl, da die einzigen Einheiten $+1$ und $-1$ sind, und wir einfach die positive Zahl wählen. 

Im Polynomring $K[x]$ hingegen sind alle konstanten Polynome ungleich null invertierbar. Wenn wir einen gemeinsamen Teiler mit einer solchen Konstanten multiplizieren, bleibt er ein größter gemeinsamer Teiler. 

Das heißt, der $\operatorname{ggT}$ ist im Polynomring nur eindeutig bis auf Multiplikation mit einer Konstanten ungleich null --- also bis auf Multiplikation mit Elementen in $K^*$. Eindeutigkeit erhalten wir, wenn wir zusätzlich fordern, dass das Polynom normiert ist, der Leitkoeffizient also gleich $1$ ist. Da hat man dann wieder eine eindeutige Wahl.
\end{spoken-clean}

\begin{math-stroke}[Eindeutigkeit des größten gemeinsamen Teilers]
\begin{remark}[Eindeutigkeit des $\operatorname{ggT}$ bis auf Einheiten]\label[remark]{rem:ggt-eindeutigkeit}
Im Gegensatz zu den ganzen Zahlen $\mathbb{Z}$, bei denen die Einheitengruppe nur aus $\{\pm 1\}$ besteht und somit eine kanonische Wahl des positiven größten gemeinsamen Teilers ermöglicht, ist die Einheitengruppe des Polynomrings $K[x]$ gegeben durch die konstanten Polynome ungleich Null:
\[
K^* = K \setminus \{0\}.
\]
Ist $h \in K[x]$ ein größter gemeinsamer Teiler von $f$ und $g$, so ist auch jedes konstante Vielfache $c \cdot h$ mit $c \in K^*$ ein größter gemeinsamer Teiler.

Der größte gemeinsame Teiler ist eindeutig bestimmt, wenn wir zusätzlich fordern, dass er \newterm{normiert} (englisch: \emph{monic}) ist, das heißt, dass sein Leitkoeffizient gleich $1$ ist.
\end{remark}
\end{math-stroke}

\subsection{Der euklidische Algorithmus für Polynome}

\begin{spoken-clean}[00:30:06 - 00:31:36]
Gut, und als weitere Bemerkung --- was ich jetzt hier nicht im Detail vorrechne ---: Man kann genau dasselbe machen wie in $\mathbb{Z}$ mit dem euklidischen Algorithmus. Auch in diesem Fall (Polynomring) kann man den euklidischen Algorithmus (\cref{thm:euclid-correctness}) anwenden, um den $\operatorname{ggT}$ sowie die Bézout-Koeffizienten $A$ und $B$ zu bestimmen. Also...
\end{spoken-clean}

\begin{math-stroke}[Euklidischer Algorithmus für Polynome]
\begin{remark}[Euklidischer Algorithmus]\label[remark]{rem:euklidischer-algorithmus-polynome}
Wie für die ganzen Zahlen $\mathbb{Z}$ gibt es auch für den Polynomring $K[x]$ den (erweiterten) euklidischen Algorithmus. Durch sukzessive Division mit Rest lässt sich für zwei Polynome $f, g \in K[x]$ der größte gemeinsame Teiler $h = \operatorname{ggT}(f, g)$ sowie die Bézout-Koeffizienten $A, B \in K[x]$ bestimmen, sodass gilt:
\[
A \cdot f + B \cdot g = h.
\]
\end{remark}
\end{math-stroke}

\subsection{Irreduzible Polynome}

\begin{spoken-clean}[00:31:36 - 00:35:30]
Das Nächste, was wir für die ganzen Zahlen hatten, waren die Primzahlen. Das können wir nun auch hier definieren. Wir definieren es zwar etwas anders, aber es stellt sich heraus, dass es im Wesentlichen dasselbe ist. 

Wir sagen: Ein Polynom $f \in K[x]$ heißt irreduzibel, falls $f \neq 0$ und nicht konstant ist, und für alle $g \in K[x]$ gilt: Falls $g$ ein Teiler von $f$ ist, dann ist $g$ konstant oder $f$ ist ein konstantes Vielfaches von $g$ --- das heißt, $f = c \cdot g$ für eine Konstante $c \in K^*$.
\end{spoken-clean}

\begin{math-stroke}[Irreduzible Polynome]
\begin{definition}[Irreduzibles Polynom]\label[definition]{def:irreduzibles-polynom}
Ein Polynom $f \in K[x]$ heißt \newterm{irreduzibel}, falls gilt:
\begin{itemize}
    \item $f \neq 0$,
    \item $f$ ist nicht konstant (das heißt, $\deg(f) \ge 1$),
    \item für alle $g \in K[x]$ gilt: falls $g \mid f$, dann ist $g$ konstant (eine Einheit in $K^*$) oder $f$ ist ein konstantes Vielfaches von $g$ (das heißt, $f = c \cdot g$ für ein $c \in K^*$).
\end{itemize}
Ein nicht-konstantes Polynom, das nicht irreduzibel ist, heißt \newterm{reduzibel}.
\end{definition}
\end{math-stroke}

\begin{meta-note}[Tafelreinigung]
Der Dozent wischt die linke und mittlere Tafel, um Platz für die nächsten Sätze über die Primfaktorzerlegung von Polynomen zu schaffen.
\end{meta-note}

\subsection{Primfaktorzerlegung in \texorpdfstring{$K[x]$}{K[x]}}

\begin{spoken-clean}[00:35:30 - 00:38:02]
Satz: Jedes Polynom ungleich null lässt sich als Produkt von irreduziblen Polynomen schreiben. Das ist genau das Analogon zum Hauptsatz der Arithmetik (\cref{thm:hauptsatz-der-arithmetik}), den wir für die ganzen Zahlen gesehen haben. 

Wenn wir ein Polynom $f \neq 0$ haben, dann existieren ein $k \ge 0$, irreduzible, normierte Polynome $p_1, \dots, p_k$ und eine Konstante $c \in K^*$, sodass $f = c \cdot p_1 \cdot p_2 \cdots p_k$ ist. Diese Darstellung ist eindeutig bis auf die Reihenfolge der Faktoren.
\end{spoken-clean}

\begin{math-stroke}[Eindeutige Primfaktorzerlegung in K[x]]
\begin{theorem}[Hauptsatz für Polynomringe]\label[theorem]{thm:hauptsatz-polynomringe}
Sei $f \in K[x] \setminus \{0\}$. Dann existiert ein $k \ge 0$, irreduzible, normierte Polynome $p_1, \dots, p_k \in K[x]$ und eine Konstante $c \in K^*$, sodass:
\[
f = c \cdot p_1 \cdot p_2 \cdots p_k.
\]
Diese Darstellung ist eindeutig bis auf die Reihenfolge der Faktoren $p_1, \dots, p_k$.
\end{theorem}
\end{math-stroke}

\subsection{Charakterisierung irreduzibler Polynome}

\begin{spoken-clean}[00:38:02 - 00:39:47]
Das Nächste, was wir haben, ist das Analogon zum Lemma von Euklid (\cref{prop:lemma-von-euklid}). Es besagt: Wenn ein irreduzibles Polynom ein Produkt teilt, dann muss es bereits mindestens einen der Faktoren teilen. 

Seien $f, g \in K[x]$ und $p$ ein nicht-konstantes Polynom. Dann sind folgende Aussagen äquivalent: Erstens, $p$ ist irreduzibel; und zweitens, falls $p \mid (f \cdot g)$, dann gilt $p \mid f$ oder $p \mid g$.
\end{spoken-clean}

\begin{math-stroke}[Lemma von Euklid für Polynome]
\begin{theorem}[Charakterisierung irreduzibler Polynome]\label[theorem]{thm:charakterisierung-irreduzibel}
Sei $p \in K[x] \setminus K$ ein nicht-konstantes Polynom. Dann sind äquivalent:
\begin{enumerate}[label=\textbf{(\alph*)}]
    \item $p$ ist irreduzibel.
    \item Für alle $f, g \in K[x]$ gilt: Falls $p \mid (f \cdot g)$, dann gilt $p \mid f$ oder $p \mid g$.
\end{enumerate}
\end{theorem}
\end{math-stroke}

\begin{ai-global-state-checkpoint-invisible-content}
% timestamp: 00:11:00
% topic: Irreduzibilität und Primfaktorzerlegung in K[x]
% board_state: def:irreduzibles-polynom, thm:hauptsatz-polynomringe, thm:charakterisierung-irreduzibel
% next_goal: Einführung von Kongruenzen modulo f in K[x]
% open_loops: none
\end{ai-global-state-checkpoint-invisible-content}

\subsection{Kongruenzen modulo eines Polynoms}

\begin{spoken-clean}[00:39:47 - 00:42:42]
Jetzt gehen wir über zu Kongruenzen modulo ein Polynom $f$. Das ist genau dasselbe, was wir für die ganzen Zahlen gemacht haben, wo wir modulo eine Zahl $n$ gerechnet haben. Wir definieren das nun für Polynome. 

Wir sagen: Zwei Polynome $a$ und $b$ sind kongruent modulo $f$ --- geschrieben als $a \equiv b \pmod f$ ---, falls $f$ die Differenz $a - b$ teilt. Das heißt, falls ein Polynom $h \in K[x]$ existiert, sodass $a - b = h \cdot f$ gilt.
\end{spoken-clean}

\begin{math-stroke}[Kongruenz modulo eines Polynoms]
\begin{definition}[Kongruenz modulo $f$]\label[definition]{def:kongruenz-polynom}
Seien $f, a, b \in K[x]$ mit $f \neq 0$. Wir sagen, $a$ und $b$ sind \newterm{kongruent modulo $f$}, geschrieben als:
\[
a \equiv b \pmod f,
\]
falls $f$ die Differenz $a - b$ teilt, das heißt, falls ein $h \in K[x]$ existiert mit:
\[
a - b = h \cdot f.
\]
\end{definition}
\end{math-stroke}

\subsection{Der Restklassenring \texorpdfstring{$K[x]/\langle f \rangle$}{K[x]/(f)}}

\begin{spoken-clean}[00:42:42 - 00:44:24]
Und genau wie bei den ganzen Zahlen können wir zeigen, dass dies eine Äquivalenzrelation ist. Das ist eine sehr einfache Übung, das zu verifizieren. 

Wir können dann die Menge der Äquivalenzklassen betrachten. Wir bezeichnen diese Menge mit $K[x]/\langle f \rangle$. Das ist die Menge der Restklassen modulo $f$. Also...
\end{spoken-clean}

\begin{math-stroke}[Äquivalenzrelation und Restklassen]
\begin{proposition}[Kongruenz als Äquivalenzrelation]\label[proposition]{prop:kongruenz-polynom-aequivalenz}
Die Relation der Kongruenz modulo $f$ ist eine Äquivalenzrelation auf $K[x]$.
\end{proposition}

\begin{short-proof}[Beweis von \cref{prop:kongruenz-polynom-aequivalenz}]
Die Reflexivität, Symmetrie und Transitivität folgen unmittelbar aus den Teilbarkeitseigenschaften im Polynomring $K[x]$ analog zum Fall der ganzen Zahlen $\mathbb{Z}$.
\end{short-proof}
\end{math-stroke}

\begin{spoken-clean}[00:44:24 - 00:48:22]
Und wir können auf dieser Menge der Restklassen eine Addition und eine Multiplikation definieren. Wir definieren die Addition zweier Klassen einfach als die Klasse der Summe und die Multiplikation als die Klasse des Produkts. 

Wir müssen natürlich zeigen, dass dies wohldefiniert ist --- das heißt, unabhängig von der Wahl der Repräsentanten. Aber das zeigt man genau gleich wie für $\mathbb{Z}$. Also...
\end{spoken-clean}

\begin{math-stroke}[Der Restklassenring K[x]/(f)]
\begin{definition}[Restklassenring]\label[definition]{def:restklassenring-polynom}
Die Menge der Äquivalenzklassen bezüglich der Kongruenz modulo $f$ wird mit $K[x]/\langle f \rangle$ (oder auch $K[x]/(f)$) bezeichnet. Die Elemente heißen \newterm{Restklassen} und werden geschrieben als:
\[
[a] := a + \langle f \rangle := \{ b \in K[x] \mid b \equiv a \pmod f \}.
\]
Wir definieren auf $K[x]/\langle f \rangle$ eine Addition und eine Multiplikation durch:
\begin{align*}
[a] + [b] &:= [a + b], \\
[a] \cdot [b] &:= [a \cdot b].
\end{align*}
\end{definition}

\begin{proposition}[Wohldefiniertheit]\label[proposition]{prop:wohldefiniertheit-operationen}
Die Addition und Multiplikation auf $K[x]/\langle f \rangle$ sind wohldefiniert, das heißt, unabhängig von der Wahl der Repräsentanten.
\end{proposition}
\end{math-stroke}

\begin{ai-global-state-checkpoint-invisible-content}
% timestamp: 00:19:00
% topic: Restklassenring K[x]/(f) und Wohldefiniertheit
% board_state: def:kongruenz-polynom, prop:kongruenz-polynom-aequivalenz, def:restklassenring-polynom
% next_goal: Repräsentantensystem mittels Division mit Rest bestimmen
% open_loops: none
\end{ai-global-state-checkpoint-invisible-content}

\subsection{Eindeutige Repräsentanten im Restklassenring}

\begin{spoken-clean}[00:48:22 - 00:53:37]
Jetzt ist die Frage: Wie können wir diese Restklassen beschreiben? Bei den ganzen Zahlen modulo $n$ hatten wir die Zahlen $0$ bis $n-1$ als Repräsentanten. Hier können wir die Division mit Rest benutzen. 

Für jedes Polynom $g$ gibt es eindeutige Polynome $q$ und $r$, sodass $g = q \cdot f + r$ mit $\deg(r) < \deg(f)$ gilt. Das heißt, jede Restklasse $[g]$ enthält genau einen Repräsentanten $r$ mit einem Grad, der strikt kleiner als der Grad von $f$ ist. Also...
\end{spoken-clean}

\begin{math-stroke}[Repräsentantensystem für K[x]/(f)]
\begin{theorem}[Eindeutige Repräsentanten]\label[theorem]{thm:eindeutige-repraesentanten}
Sei $f \in K[x]$ mit $\deg(f) = n \ge 1$. Jede Restklasse $[g] \in K[x]/\langle f \rangle$ besitzt genau einen Repräsentanten $r \in K[x]$ mit:
\[
\deg(r) < \deg(f).
\]
Das heißt, die Abbildung:
\[
\varphi \colon \{ r \in K[x] \mid \deg(r) < n \} \to K[x]/\langle f \rangle, \quad r \mapsto [r]
\]
ist eine BJektion.
\end{theorem}

\begin{short-proof}[Beweis von \cref{thm:eindeutige-repraesentanten}]
\textbf{Surjektivität (Existenz):} Nach dem Satz über die Division mit Rest (\cref{thm:division-mit-rest}) existieren für jedes $g \in K[x]$ Polynome $q, r \in K[x]$ mit:
\[
g = q \cdot f + r \quad \text{und} \quad \deg(r) < \deg(f).
\]
Daraus folgt $g - r = q \cdot f \implies g \equiv r \pmod f$, also $[g] = [r]$.

\textbf{Injektivität (Eindeutigkeit):} Seien $r_1, r_2 \in K[x]$ mit $\deg(r_1), \deg(r_2) < \deg(f)$ und $[r_1] = [r_2]$. Dann gilt:
\[
r_1 \equiv r_2 \pmod f \implies f \mid (r_1 - r_2).
\]
Angenommen, $r_1 - r_2 \neq 0$. Dann gilt nach der Gradformel [Anmerkung: $r_1 - r_2 = h \cdot f$ für ein $h \neq 0$, woraus $\deg(r_1 - r_2) = \deg(h) + \deg(f) \ge \deg(f)$ folgt]:
\[
\deg(f) \le \deg(r_1 - r_2).
\]
Andererseits gilt jedoch:
\[
\deg(r_1 - r_2) \le \max\{\deg(r_1), \deg(r_2)\} < \deg(f),
\]
was ein Widerspruch ist. Folglich muss $r_1 - r_2 = 0 \implies r_1 = r_2$ gelten.
\end{short-proof}
\end{math-stroke}

\subsection{Die Körper-Eigenschaft von \texorpdfstring{$K[x]/\langle f \rangle$}{K[x]/(f)}}

\begin{spoken-clean}[00:53:37 - 00:57:19]
Und jetzt kommt der ganz wichtige Satz, der uns sagt, wann dieser Restklassenring ein Körper ist. Bei den ganzen Zahlen modulo $n$ war das genau dann ein Körper, wenn $n$ eine Primzahl war. Hier ist es genau dann ein Körper, wenn das Polynom $f$ irreduzibel ist. Das ist eine wunderschöne Analogie! Also...
\end{spoken-clean}

\begin{math-stroke}[Wann ist K[x]/(f) ein Körper?]
\begin{theorem}[Körper-Eigenschaft des Restklassenrings]\label[theorem]{thm:koerper-eigenschaft}
Sei $f \in K[x] \setminus K$ ein nicht-konstantes Polynom. Der Restklassenring $K[x]/\langle f \rangle$ ist genau dann ein Körper, wenn $f$ irreduzibel ist.
\end{theorem}

\begin{short-proof}[Beweis von \cref{thm:koerper-eigenschaft}]
\textbf{Hinrichtung ($\implies$):} Angenommen, $f$ ist reduzibel. Dann existieren nicht-konstante Polynome $g, h \in K[x]$ mit $f = g \cdot h$ und $\deg(g), \deg(h) < \deg(f)$.
Im Restklassenring $K[x]/\langle f \rangle$ gilt dann:
\[
[g] \cdot [h] = [g \cdot h] = [f] = [0].
\]
Da $\deg(g), \deg(h) < \deg(f)$, gilt $[g] \neq [0]$ und $[h] \neq [0]$. Somit besitzt $K[x]/\langle f \rangle$ Nullteiler und kann kein Körper sein.

\textbf{Rückrichtung ($\impliedby$):} Sei $f$ irreduzibel und $[a] \in K[x]/\langle f \rangle$ mit $[a] \neq [0]$, das heißt, $f \nmid a$.
Da $f$ irreduzibel ist, sind die einzigen Teiler von $f$ die Einheiten und die konstanten Vielfachen von $f$ selbst. Da $f \nmid a$, muss der größte gemeinsame Teiler von $a$ und $f$ eine Einheit sein, also $\operatorname{ggT}(a, f) = 1$.
Nach dem Satz von Bézout (\cref{thm:ggt-bezout}) existieren Polynome $A, B \in K[x]$ mit:
\[
A \cdot a + B \cdot f = 1.
\]
Gehen wir zu den Restklassen über, erhalten wir:
\[
[A \cdot a + B \cdot f] = [1] \implies [A] \cdot [a] + [B] \cdot [0] = [1] \implies [A] \cdot [a] = [1].
\]
Somit ist $[A]$ das multiplikative Inverse von $[a]$ in $K[x]/\langle f \rangle$. Da jedes Element ungleich Null ein Inverses besitzt, ist $K[x]/\langle f \rangle$ ein Körper.
\end{short-proof}

\begin{explanation-of-steps}[Konstruktion endlicher Körper (Galois-Körper)]
Dieses Theorem bildet das Fundament für die Konstruktion aller endlichen Körper:
\begin{itemize}
    \item \textbf{Analogie $\mathbb{Z}$ vs. $K[x]$:} Ein irreduzibles Polynom $f \in K[x]$ spielt im Polynomring genau dieselbe Rolle wie eine Primzahl $p \in \mathbb{Z}$ im Ring der ganzen Zahlen.
    \item \textbf{Erzeugung endlicher Körper $\mathbb{F}_{p^k}$:} Wählen wir $K = \mathbb{F}_p = \mathbb{Z}/p\mathbb{Z}$ als Grundkörper und ein irreduzibles Polynom $f \in \mathbb{F}_p[x]$ vom Grad $k$, so ist der Quotientenring $\mathbb{F}_p[x]/\langle f \rangle$ ein Körper mit exakt $p^k$ Elementen.
\end{itemize}
\end{explanation-of-steps}
\end{math-stroke}

\begin{ai-global-state-checkpoint-invisible-content}
% timestamp: 00:28:00
% topic: Körper-Eigenschaft von K[x]/(f)
% board_state: thm:koerper-eigenschaft, thm:eindeutige-repraesentanten
% next_goal: Konstruktion eines konkreten endlichen Körpers mit 4 Elementen
% open_loops: none
\end{ai-global-state-checkpoint-invisible-content}

\subsection{Konstruktion und Eigenschaften des endlichen Körpers \texorpdfstring{$\mathbb{F}_4$}{F4}}

\begin{spoken-clean}[00:57:19 - 00:58:57]
Schauen wir uns ein konkretes Beispiel an: Wir nehmen den Körper mit zwei Elementen, $\mathbb{F}_2$, und das Polynom $f(x) = x^2 + x + 1$. Das ist ein Polynom vom Grad $2$. Wir können leicht prüfen, ob es irreduzibel ist. Da es Grad $2$ hat, ist es genau dann reduzibel, wenn es eine Nullstelle in $\mathbb{F}_2$ hat. Aber $f(0) = 1 \neq 0$ und $f(1) = 1 \neq 0$. Also hat es keine Nullstelle und ist irreduzibel.

Wir arbeiten also über dem Körper mit zwei Elementen. Dass dieses Polynom irreduzibel ist, können Sie als kleine Übung selbst zeigen. Das ist nicht schwierig: Da es Grad $2$ hat, müsste es, wenn es nicht irreduzibel wäre, ein Produkt von zwei Polynomen von Grad $1$ sein. Wir kennen die Polynome von Grad $1$ über $\mathbb{F}_2$ --- da gibt es nicht sehr viele, die kann man alle einfach durchmultiplizieren.

Man kann auch einfach prüfen, dass es keine Nullstelle in $\mathbb{F}_2$ hat, das ist also sehr einfach. Für Polynome von hohem Grad ist es manchmal schwieriger zu zeigen, dass sie irreduzibel sind, aber hier ist es extrem einfach.
\end{spoken-clean}

\begin{spoken-clean}[00:58:57 - 01:00:12]
Das heißt, wir wissen, dass $\mathbb{F}_2[x]/\langle f \rangle$ ein Körper ist. Wie viele Elemente hat dieser Körper? Wir können sie sogar alle hinschreiben. Wie wir gesehen haben, werden die Elemente hier durch Polynome von Grad strikt kleiner als $2$ repräsentiert --- das heißt, durch Polynome vom Grad $0$ und $1$. Davon gibt es genau vier. Das heißt, wir haben einen Körper mit vier Elementen.

Nämlich explizit: Wir haben die Restklassen $[0]$, $[1]$, $[x]$ und $[x+1]$.
\end{spoken-clean}

\begin{math-stroke}[Konstruktion eines endlichen Körpers mit 4 Elementen]
\begin{example}[Der Körper $\mathbb{F}_4$]\label[example]{ex:koerper-f4}
Wir betrachten den Grundkörper $K = \mathbb{F}_2 = \{0, 1\}$ und das Polynom:
\[
f(x) = x^2 + x + 1 \in \mathbb{F}_2[x].
\]

\textbf{1. Irreduzibilität von $f$:} \\
Da $\deg(f) = 2$, ist $f$ genau dann reduzibel, wenn es einen Linearfaktor besitzt, also eine Nullstelle in $\mathbb{F}_2$ hat. Wir prüfen die beiden Elemente von $\mathbb{F}_2$:
\begin{align*}
f(0) &= 0^2 + 0 + 1 = 1 \neq 0, \\
f(1) &= 1^2 + 1 + 1 = 1 \neq 0 \quad (\text{in } \mathbb{F}_2).
\end{align*}
Da $f$ keine Nullstelle in $\mathbb{F}_2$ besitzt, ist $f$ irreduzibel. Nach \cref{thm:koerper-eigenschaft} ist der Restklassenring:
\[
\mathbb{F}_4 := \mathbb{F}_2[x]/\langle x^2 + x + 1 \rangle
\]
ein Körper.

\textbf{2. Elemente des Körpers:} \\
Da $\deg(f) = 2$, besteht das eindeutige Repräsentantensystem nach \cref{thm:eindeutige-repraesentanten} aus allen Polynomen vom Grad $< 2$:
\[
\mathbb{F}_4 = \{ [0], [1], [x], [x+1] \}.
\]
Dieser Körper besitzt somit genau $2^2 = 4$ Elemente.

\textbf{3. Multiplikationstabelle (Beispiele):} \\
Wir berechnen beispielhaft einige Produkte im Körper $\mathbb{F}_4$ unter Ausnutzung der Relation $[x^2 + x + 1] = [0] \implies [x^2] = [x + 1]$ (da $-1 = 1$ in $\mathbb{F}_2$):
\begin{itemize}
    \item \textbf{Produkt $[x] \cdot [x]$:}
    \[
    [x] \cdot [x] = [x^2] = [x + 1].
    \]
    \item \textbf{Produkt $[x] \cdot [x+1]$:}
    \begin{align*}
    [x] \cdot [x+1] &= [x^2 + x] = [x^2] + [x] \\
    &= [x+1] + [x] = [2x + 1] = [1] \quad (\text{da } 2 = 0 \text{ in } \mathbb{F}_2).
    \end{align*}
    Somit sind $[x]$ und $[x+1]$ multiplikativ invers zueinander: $[x]^{-1} = [x+1]$.
\end{itemize}
\end{example}
\end{math-stroke}

\begin{spoken-clean}[01:00:12 - 01:02:59]
Man muss hier noch schauen, wie man diese miteinander multipliziert. Mit $[0]$ und $[1]$ ist das natürlich klar. Für die anderen beiden können wir das berechnen: Wenn wir $[x] \cdot [x+1]$ multiplizieren, erhalten wir die Klasse von $x^2 + x$. Da aber $x^2 + x + 1 \equiv 0 \pmod f$ gilt, ist das dasselbe wie die Klasse von $[-1]$. Und da $-1 = 1$ in $\mathbb{F}_2$ ist, ist das dasselbe wie die Klasse von $[1]$.

Das heißt, $[x]$ und $[x+1]$ sind multiplikativ invers zueinander. Dann können wir noch schauen, was das Quadrat von $[x]$ ist: $[x]^2 = [x^2] = [x+1]$, genau! So kann man wunderbar damit rechnen. Das heißt: Wenn man ein irreduzibles Polynom hat, erhält man einen Körper.
\end{spoken-clean}

\begin{math-stroke}[Multiplikation in \texorpdfstring{$\mathbb{F}_4$}{F4}]
Im Körper $\mathbb{F}_4$ gelten folgende Multiplikationsregeln:
\begin{align*}
[x] \cdot [x+1] &= [x^2 + x] = [1], \\
[x]^2 &= [x+1].
\end{align*}
\end{math-stroke}

\begin{spoken-clean}[01:02:59 - 01:05:39]
Und vielleicht noch ein Beispiel, das Sie wahrscheinlich schon kennen --- das ist eine gute Übung, das wirklich zu zeigen ---: Der Körper der komplexen Zahlen $\mathbb{C}$ ist isomorph zum Quotienten des Polynomrings $\mathbb{R}[x]$ modulo dem irreduziblen Polynom $x^2 + 1$. Das ist eine schöne Übung.
\end{spoken-clean}

\begin{math-stroke}[Isomorphie der komplexen Zahlen]
\begin{example}[Die komplexen Zahlen als Restklassenkörper]\label[example]{ex:komplexe-zahlen-restklassen}
Der Körper der komplexen Zahlen $\mathbb{C}$ ist isomorph zum Restklassenkörper des Polynomrings über den reellen Zahlen $\mathbb{R}[x]$ modulo des irreduziblen Polynoms $x^2 + 1$:
\[
\mathbb{C} \cong \mathbb{R}[x]/\langle x^2 + 1 \rangle \quad (\text{Übung}).
\]
\end{example}
\end{math-stroke}

\subsection{Klassifikation endlicher Körper}

\begin{meta-note}[Tafelreinigung]
Der Dozent wischt die Tafel, um Platz für den Klassifikationssatz endlicher Körper zu machen.
\end{meta-note}

\begin{spoken-clean}[01:05:39 - 01:08:19]
Was uns jetzt interessiert, ist der Fall, dass $K$ ein endlicher Körper ist --- zum Beispiel $K = \mathbb{F}_p$. Dann gibt es für jedes irreduzible Polynom $f \in \mathbb{F}_p[x]$ vom Grad $n$ einen Körper $\mathbb{F}_p[x]/\langle f \rangle$, der, wie wir gesehen haben, genau $p^n$ Elemente besitzt. Und wir haben gesehen: Es gibt unendlich viele irreduzible Polynome von beliebig hohem Grad in $\mathbb{F}_p[x]$.

Der folgende Satz --- ich schreibe ihn nur hin und verzichte auf den Beweis --- besagt: Zu jeder Primzahl $p$ und jedem $n \ge 1$ existiert ein Körper der Ordnung $p^n$. Mehr noch: Er existiert nicht nur, sondern ist auch eindeutig bis auf Isomorphie. Jeder endliche...
\end{spoken-clean}

\begin{math-stroke}[Klassifikationssatz endlicher Körper]
\begin{theorem}[Existenz und Eindeutigkeit endlicher Körper]\label[theorem]{thm:klassifikation-endliche-koerper}
Zu jeder Primzahl $p$ und jeder natürlichen Zahl $n \ge 1$ existiert ein Körper der Ordnung $p^n$, welcher bis auf Isomorphie eindeutig bestimmt ist.

Jeder endliche Körper ist von dieser Form.
\end{theorem}

\begin{remark}[Freiwilliges Nachlesen]\label[remark]{rem:freiwillig-nachlesen}
Der Beweis dieses Satzes wird in der Vorlesung Algebra I/II behandelt und kann bei Interesse selbstständig nachgelesen werden.
\end{remark}
\end{math-stroke}

\begin{spoken-clean}[01:08:19 - 01:11:09]
Und jeder endliche Körper ist von dieser Form. Jeder endliche Körper hat also $p^n$ Elemente für eine Primzahl $p$ und ein $n \ge 1$. Bis auf Isomorphismus gibt es genau einen Körper von dieser Ordnung. Das heißt, endliche Körper sind vollständig klassifiziert!

Was Sie mit Ihrem jetzigen Wissen bereits gut sehen können --- und auch gerne nachlesen dürfen ---, ist, dass jeder endliche Körper die Ordnung $p^n$ haben muss. Das ist nicht so schwierig zu sehen: Sie müssen nur bemerken, dass ein endlicher Körper ein Vektorraum über seinem Primkörper $\mathbb{F}_p$ ist, und daraus folgt das direkt.

Um zu sehen, dass es für jedes $n$ einen Körper der Ordnung $p^n$ gibt, muss man ein bisschen arbeiten, aber auch das ist machbar. Etwas schwieriger ist es zu zeigen, dass alle Körper der Ordnung $p^n$ isomorph sind; da muss man noch ein bisschen mehr arbeiten. Sie werden das auf jeden Fall in der Algebra-Vorlesung sehen --- oder Sie können es freiwillig nachlesen.
\end{spoken-clean}

\begin{math-stroke}[Struktur endlicher Körper]
Jeder endliche Körper $E$ besitzt notwendigerweise eine Primzahlcharakteristik $p>0$ und enthält somit einen zu $\mathbb{F}_p$ isomorphen Teilkörper (den Primkörper). Man kann $E$ daher als Vektorraum über $\mathbb{F}_p$ auffassen. Da $E$ endlich ist, ist die Dimension $n = \dim_{\mathbb{F}_p}(E)$ endlich. 
\begin{explanation-of-steps}[Kardinalität des Vektorraums]
Da jeder $n$-dimensionale Vektorraum über einem Körper $K$ isomorph zu $K^n$ ist, gibt es für jede der $n$ Basisrichtungen genau $|K| = p$ Wahlmöglichkeiten. Daraus ergibt sich die Gesamtzahl der Elemente in $E$:
\end{explanation-of-steps}
\[
|E| = p^n.
\]
\end{math-stroke}

\begin{spoken-clean}[01:11:09 - 01:12:58]
Was vielleicht noch wichtig ist und was Sie auf jeden Fall in Ihr mathematisches Leben mitnehmen sollten, ist: Nicht alle endlichen Körper sind von der Form $\mathbb{F}_p$. Unser Beispiel mit vier Elementen ist kein Körper der Form $\mathbb{F}_p$. Das ist der Körper $\mathbb{F}_4$ (d.\,h. $\mathbb{F}_{2^2}$), aber das ist eben nicht der Restklassenring $\mathbb{Z}/4\mathbb{Z}$. Das ist extrem wichtig zu sehen!
\end{spoken-clean}

\begin{math-stroke}[Unterscheidung von Körpern und Ringen]
\begin{remark}[Strukturelle Unterschiede]\label[remark]{rem:strukturelle-unterschiede}
Es ist essenziell, den endlichen Körper $\mathbb{F}_4$ (welcher isomorph zu $\mathbb{F}_2[x]/\langle x^2+x+1\rangle$ ist) nicht mit dem Restklassenring $\mathbb{Z}/4\mathbb{Z}$ zu verwechseln. Letzterer besitzt Nullteiler (da $[2] \cdot [2] = [0]$) und ist somit kein Körper.
\end{remark}
\end{math-stroke}

\section{Prüfungsinformationen und Ausblick}

\begin{meta-note}[Tafelreinigung und Übergang]
Der Dozent wischt die Tafel vollständig, um organisatorische Details zur bevorstehenden Prüfung zu notieren.
\end{meta-note}

\subsection{Prüfungsinhalt und Übungen}

\begin{spoken-clean}[01:12:58 - 01:16:29]
Gut, so viel erst einmal als kleiner Teaser für die Algebra und um zu sehen, was man da alles machen kann. Gibt es dazu noch Fragen? Zu den Körpern oder Ähnlichem? Nein?

Dann möchte ich gerne die verbleibende Zeit nutzen, um noch kurz mit Ihnen über die Prüfung und die Prüfungsvorbereitung zu sprechen. Die erste Frage betrifft meistens den Prüfungsinhalt: Prüfungsrelevant ist der gesamte Inhalt der Vorlesung --- also alles, was wir in der Vorlesung und in den Übungen behandelt haben.
\end{spoken-clean}

\begin{spoken-clean}[01:16:29 - 01:18:37]
Das beinhaltet natürlich nicht die Teile im Skript, die wir nicht behandelt haben. Ich werde heute Abend noch ein Dokument hochladen, in dem die behandelten Kapitel genau aufgelistet sind. In den letzten drei Wochen sind wir ja eher quer durch die letzten drei Kapitel gegangen und haben einige Dinge etwas anders gemacht. Da werde ich die Notizen der letzten drei Wochen hochladen.

Schreibt hier zufällig jemand die Vorlesung in LaTeX mit? Das gibt es ja immer wieder. Wenn Sie möchten --- Sie müssen natürlich nicht ---, können Sie mir gerne Ihre Mitschrift der letzten drei Wochen schicken. Dann kann ich das schnell überfliegen und allen zugänglich machen. Dann haben alle eine saubere Grundlage, die genau dem prüfungsrelevanten Stoff entspricht. Wäre das gut? Ich würde es dann kurz durchlesen und eventuell noch ein bisschen anpassen. Cool, vielen Dank, super! Ich glaube, alle hier sind Ihnen dankbar, wenn Sie das machen.
\end{spoken-clean}

\begin{spoken-clean}[01:18:37 - 01:20:27]
Und sonst müsste ich die Videos irgendwie bei einem KI-Modell hochladen und hoffen, dass da etwas Vernünftiges herauskommt --- aber das klappt meistens nicht so gut. Das ist also der Vorlesungsinhalt. Ein ganz wichtiger Punkt ist allerdings: Auch die Übungen sind Teil des Prüfungsstoffs. Die Übungsblätter sind, gerade in dieser Vorlesung, ein essenzieller Bestandteil. Es ist wichtig, dass Sie unter dem Semester an diesen Aufgaben gearbeitet haben.

Wir haben in der Vorlesung vielleicht etwas weniger Stoff behandelt, weil ich nicht zu viel hineinquetschen wollte; dafür sind die Übungsblätter eventuell etwas umfangreicher als in anderen Vorlesungen. Ich hoffe, Sie haben diese unter dem Semester aktiv bearbeitet, diskutiert und verstanden. Es gibt ja auch Musterlösungen zu allen Blättern. Aber wie gesagt: Die Übungen sind fester Bestandteil des Stoffs --- alles, was dort eingeführt wurde, ist absolut prüfungsrelevant. Es gibt auch diese Woche noch ein Übungsblatt $14$. Wir haben also wirklich jede Woche ein Blatt gehabt. Aber das waren ja hoffentlich meistens lustige oder zumindest interessante Aufgaben.
\end{spoken-clean}

\begin{nice-box}[Prüfungsstoff und Relevanz]
\begin{itemize}
    \item \textbf{Vorlesungsinhalt:} Alle behandelten Kapitel des Skripts (eine detaillierte Liste wird online bereitgestellt).
    \item \textbf{Übungsblätter:} Die wöchentlichen Übungsblätter (inklusive des aktuellen Blatts $14$) sind ein integraler Bestandteil des Prüfungsstoffs.
    \item \textbf{Zusatzmaterial:} Notizen zu den letzten drei Wochen der Vorlesung werden hochgeladen.
\end{itemize}
\end{nice-box}

\subsection{Prüfungsformat und Hilfsmittel}

\begin{spoken-clean}[01:20:27 - 01:22:57]
Kommen wir zur Prüfungsform: Die Prüfung dauert zwei Stunden und findet bereits in der ersten Prüfungswoche statt. Zwei Stunden sind gar nicht so lang, das geht wie im Flug vorbei. Etwas ein Drittel der Prüfung --- das ist kein exakter Wert, sondern ein Richtwert --- wird aus Multiple-Choice-Fragen bestehen. Das sind entweder Single-Choice-Fragen, bei denen Sie eine Antwort aus vier Optionen auswählen, oder Richtig-Falsch-Fragen. Es gibt keine Minuspunkte für falsche Antworten. Das macht man heute aus verschiedenen Gründen eigentlich nicht mehr. Das heißt für Sie: Kreuzen Sie auf jeden Fall etwas an, selbst wenn Sie die Antwort nicht wissen! Es ist zwar didaktisch unschön, dass der Zufall eine Rolle spielt, aber so sind die Regeln.

Die restlichen zwei Drittel der Prüfung bestehen aus offenen Aufgaben. Hier spielt die Art und Weise, wie Sie Ihre Antworten aufschreiben, eine ganz entscheidende Rolle --- besonders in dieser Vorlesung, in der es darum geht zu lernen, wie man Mathematik präzise und elegant formuliert. Geben Sie sich also Mühe, wirklich schöne, saubere Argumente aufzuschreiben. Wie immer gilt: Bewertet wird ausschließlich das, was auf dem Papier steht, nicht das, was Sie im Kopf hatten. Manchmal ist das den Studierenden nicht ganz klar. Da heißt es im Nachhinein oft: \qt{Ja, aber ich habe eigentlich das gemeint, und das sollte doch aus dem Kontext klar sein!} Nein, bewertet wird nur das, was Sie tatsächlich aufschreiben.

Das haben Sie aber hoffentlich das ganze Semester über fleißig geübt, indem Sie Ihre Übungen abgegeben, Feedback erhalten und versucht haben, Ihre Argumentation zu verbessern. Machen Sie es bitte nicht so wie manche Studierende im ersten Jahr, die einfach eine Reihe von Behauptungen hinschreiben und zur Sicherheit auf jeder Zeile einen Doppelimplikationspfeil ($\iff$) ergänzen, selbst wenn dieser in der Hälfte der Fälle logisch gar nicht stimmt. Seien Sie vorsichtig, argumentieren Sie sauber!
\end{spoken-clean}

\begin{spoken-clean}[01:22:57 - 01:23:49]
Sie dürfen ein A4-Blatt --- also zwei Seiten --- Zusammenfassung, Formelsammlung oder was auch immer mitnehmen. Es gibt keinerlei Einschränkungen, was darauf stehen darf. Sie können auch ein Gedicht oder ein Foto von... ja, völlig egal, was Sie wollen. Es wäre lächerlich zu verlangen, dass es handgeschrieben sein muss, und wir können auch nicht überprüfen, ob Sie es selbst verfasst haben. Bringen Sie einfach ein A4-Blatt mit, das doppelseitig bedruckt oder beschrieben ist. Was darauf steht, ist uns völlig egal.
\end{spoken-clean}

\begin{math-stroke}[Struktur der Prüfung]
\begin{itemize}
    \item \textbf{Dauer:} $2$ Stunden.
    \item \textbf{Format:}
    \begin{itemize}
        \item \textbf{$\approx 1/3$} Multiple-Choice-Fragen (keine negativen Punkte für falsche Antworten).
        \item \textbf{$\approx 2/3$} offene Aufgaben (Bewertung der mathematischen Präzision und Argumentation).
    \end{itemize}
    \item \textbf{Erlaubte Hilfsmittel:} \textbf{$1$} A4-Blatt ($2$ Seiten) Zusammenfassung (beliebiger Inhalt, gedruckt oder handgeschrieben).
\end{itemize}
\end{math-stroke}

\begin{student-interaction}[Studentenfrage]
Müssen wir die ganzen Axiome auswendig lernen, wie z.B. die Peano-Axiome, oder stehen die auf der Prüfung drauf?
\end{student-interaction}

\begin{spoken-clean}[01:23:49 - 01:24:19]
Nein, ich kann das gerne noch präzisieren: Die Peano-Axiome sollten Sie vielleicht schon können, aber die ganzen logischen Axiome oder die Zermelo-Fraenkel-Axiome müssen Sie jetzt nicht alle auswendig lernen. Sie müssen diese Axiome verstanden haben, aber Sie müssen sie nicht auswendig aufsagen können. Wenn Sie sie anwenden sollen, würden wir sie vorgeben --- zum Beispiel in dem Stil: \qt{Hier sind fünf Axiome, verwenden Sie diese, um die folgende Aussage formal zu beweisen.} So etwas in der Art könnte kommen.
\end{spoken-clean}

\subsection{Prüfungsvorbereitung und Tipps}

\begin{spoken-clean}[01:24:19 - 01:26:29]
Es ist immer schwierig zu sagen. Ich versuche manchmal, das so zu erklären: Es gibt zwei Extreme bei Prüfungen. Auf der einen Seite haben wir die Mathe-Olympiade: Da versucht man natürlich, möglichst unerwartete, trickreiche und schwierige Fragen auszutüfteln, die niemand vorher kennt. Das andere Extrem ist die Fahrprüfung: Diese ist im Wesentlichen sehr vorhersehbar. Man weiß genau, was man können muss, und muss das dann eben abrufen.

Obwohl es auch bei der Fahrprüfung unerwartete Situationen gibt --- Sie fahren auf der Autobahn, haben das oft geübt, und genau in der Prüfung kommt dieser riesige Lastwagen und Sie müssen schauen, wie Sie trotzdem richtig reagieren. Also ist auch die Fahrprüfung nicht völlig ohne Überraschungen.

Die Frage ist nun: Wo liegen die Mathe-Prüfungen an der ETH auf dieser Skala? Natürlich nicht ganz bei der Fahrprüfung, aber sie sollten auch nicht wie die Mathe-Olympiade sein. Wir versuchen, uns irgendwo in der Mitte einzupendeln. Es ist natürlich auch für uns nicht so einfach, gute Prüfungen zu entwerfen.
\end{spoken-clean}

\begin{math-stroke}[Die Prüfungsskala]
\begin{center}
\begin{tikzpicture}[scale=1.5, >=Stealth]
    % \begin{ai-tikz-planner-invisible-content}
    % 1. Background: Horizontal line representing the scale.
    % 2. Midground: Ticks at the ends and the middle.
    % 3. Foreground: Labels "Mathe-Olympiade" (left), "Fahrprüfung" (right), and the target region (middle).
    % \end{ai-tikz-planner-invisible-content}
    \draw[thick, <->] (-3,0) -- (3,0);
    \draw[thick] (-2.5, 0.1) -- (-2.5, -0.1) node[below=0.1cm, align=center, font=\small] {Mathe-Olympiade\\(Sehr trickreich)};
    \draw[thick] (2.5, 0.1) -- (2.5, -0.1) node[below=0.1cm, align=center, font=\small] {Fahrprüfung\\(Sehr vorhersehbar)};
    
    % Target area with beautiful fill
    \fill[BrickRed!10, draw=BrickRed, dashed, thick] (0, -0.3) rectangle (1, 0.3);
    \node[BrickRed, above=0.3cm, font=\bfseries\small] at (0.5, 0) {Zielbereich};
\end{tikzpicture}
\end{center}
\end{math-stroke}

\begin{spoken-clean}[01:26:29 - 01:26:54]
Ich versuche immer, ein bisschen Mut zu machen, dass man keine allzu große Angst vor der Prüfung haben muss. Wenn man gut gelernt und die Übungen gemacht hat, dann schafft man das auch. Aber das wird manchmal fälschlicherweise so verstanden, als ob die Prüfung einfach wird. Es ist und bleibt eine Mathe-Prüfung an der ETH im ersten Jahr --- und das ist naturgemäß nicht einfach. Bereiten Sie sich also gründlich vor! Ein Kreditpunkt entspricht laut Richtlinien etwa $30$ Stunden Arbeit. Sie sind zwar im Vergleich zur Allgemeinbevölkerung alle mathematisch überdurchschnittlich begabt, aber innerhalb Ihrer Kohorte an der ETH sind die meisten von Ihnen im Durchschnitt. Daher sind $30$ Stunden Arbeit pro Credit ein sehr realistischer Richtwert, selbst wenn Sie an der Mittelschule viel weniger für Mathematik tun mussten. Schauen Sie einfach mal, wie viel Sie bisher investiert haben, um ein Gefühl dafür zu bekommen, wie viel Vorbereitung noch nötig ist.

Ich werde ein paar Prüfungen der letzten Jahre hochladen und mich bei der Erstellung der neuen Prüfung stark an diesen orientieren --- natürlich unter Ausschluss der Themen, die wir dieses Semester nicht behandelt haben. Viele von Ihnen wünschen sich eine Probeprüfung, daher werde ich eine erstellen. Ich muss aber betonen, dass die Probeprüfung nur mäßig nützlich sein wird. Das Einzige, was ich Ihnen garantieren kann, ist, dass die echte Prüfung nicht so aussehen wird wie die Probeprüfung --- alles andere wäre ja witzlos. Wenn ich eine Frage für die echte Prüfung im Kopf habe, darf sie natürlich nicht auf der Probeprüfung stehen. Erwarten Sie also nicht, dass die echte Prüfung ein Klon der Probeprüfung sein wird. Ich erstelle sie hauptsächlich zur Beruhigung, aber orientieren Sie sich beim Lernen lieber an den echten Prüfungen der Vorjahre.

Wie bereitet man sich am besten vor? Für viele von Ihnen ist es die erste Prüfungssession an der ETH. Eine sehr effiziente Methode ist es, sich selbst Prüfungsfragen auszudenken. Überlegen Sie sich: Wenn ich diese Prüfung entwerfen müsste, was würde ich fragen? Schreiben Sie diese Fragen auf. Im Gegensatz zur Probeprüfung ist hier die Wahrscheinlichkeit positiv, dass eine Ihrer Fragen tatsächlich in der echten Prüfung vorkommt! Versuchen Sie dann, diese selbst entworfenen Fragen zu lösen, oder tauschen Sie sie mit Ihren Kolleginnen und Kollegen aus. Was hingegen gar nicht funktioniert, ist, einfach nur das Skript durchzublättern, ab und zu zu nicken und zu denken: \qt{Ah ja, das habe ich verstanden.} Es ist ein riesiger Unterschied, ob man eine Argumentation im Skript nachvollziehen kann oder ob man vor einem leeren Blatt sitzt und sie selbstständig reproduzieren muss. Das ist wie beim Autofahren: Dem Fahrlehrer zuzuschauen ist das eine, selbst am Steuer zu sitzen etwas ganz anderes.

Heutzutage bereitet man sich wahrscheinlich auch viel mit generativer KI auf Prüfungen vor. Das kann durchaus hilfreich sein, man muss aber den richtigen Umgang damit finden. Verschiedene Studien zeigen, dass das Lernen mit KI oft zu einer systematischen Selbstüberschätzung führt. Studierende, die sich primär mit KI vorbereiten, schätzen ihre eigenen Fähigkeiten im Nachhinein oft zu hoch ein. Man hat das Gefühl, alles verstanden zu haben, und die KI bestätigt einem vielleicht noch, dass man bestens vorbereitet ist. Die Gefahr der Fehleinschätzung ist hier laut Studien sehr real. Seien Sie also kritisch!

Ganz wichtig ist auch: Machen Sie Pausen! Versuchen Sie nicht, zwölf Stunden am Tag zu lernen --- das ist physiologisch überhaupt nicht produktiv. Oft ist weniger mehr. Wir haben noch ein paar Minuten: Gibt es noch Fragen zur Prüfung oder zum Inhalt? Nein? Dann möchte ich gerne... Letzte Frage, Wortmeldung?
\end{spoken-clean}

\begin{math-stroke}[Prüfungskriterien und Vorbereitung]
\textbf{Prüfungskriterien:}
\begin{itemize}
    \item \textbf{Breite:} Ausgewogene Abdeckung des Vorlesungsstoffs und der Übungen.
    \item \textbf{Schwierigkeit:} Eine gesunde Mischung aus direkt lösbaren Aufgaben und anspruchsvolleren Transferfragen.
    \item \textbf{Gewichtung:} Orientierung an der Relevanz der Themen im Semester.
\end{itemize}
\end{math-stroke}

\begin{nice-box}[Tipps zur Prüfungsvorbereitung]
\begin{itemize}
    \item \textbf{Aktives Lernen:} Das Skript nicht nur passiv lesen, sondern versuchen, Sätze und Beweise auf einem leeren Blatt Papier eigenständig zu reproduzieren.
    \item \textbf{Eigene Fragen entwerfen:} Eine hervorragende Methode ist es, sich selbst Prüfungsaufgaben auszudenken und diese zu lösen.
    \item \textbf{Umgang mit KI-Tools:} Generative KI kann beim Lernen unterstützen, birgt jedoch die Gefahr, das eigene Verständnis zu überschätzen. Nutzen Sie KI-Feedback kritisch.
    \item \textbf{Pausen einlegen:} Physiologisch ist es unmöglich, $12$ Stunden am Stück produktiv zu lernen. Planen Sie feste Erholungsphasen ein.
\end{itemize}
\end{nice-box}

\begin{student-interaction}[Studentenfrage]
Wie sieht es mit der Gewichtung der...
\end{student-interaction}

\begin{spoken-clean}[01:26:54 - 01:27:19]
Ah ja, das ist ein wichtiger Punkt: Folgefehler! Ja, das ist extrem wichtig. Wenn man gleich am Anfang die Definition von etwas vergisst, ist man sonst für ein Drittel (i.e., $\approx 1/3$) der Aufgabe disqualifiziert. Ja, ein sehr guter Punkt! Mhm.
\end{spoken-clean}

\begin{student-interaction}[Studentenfrage]
Also allgemein, dass die Prüfung so ist...
\end{student-interaction}

\begin{spoken-clean}[01:27:19 - 01:27:49]
Ja, das versucht man natürlich zu vermeiden. Wir achten darauf, dass sich die Prüfung nicht zu stark auf ein einzelnes Thema konzentriert. Das betrifft auch den Schwierigkeitsgrad: Wir wollen Sie nicht übermäßig ermüden. Es ist zwar schöner, wenn Sie sich zwanzig Minuten lang tief in ein Thema hineindenken können, aber wir müssen in der Prüfung trotzdem differenzieren. Ja, noch eine letzte Wortmeldung?
\end{spoken-clean}

\begin{student-interaction}[Studentenfrage]
Eher so eine Frage, und zwar, was soll...
\end{student-interaction}

\begin{spoken-clean}[01:27:49 - 01:28:29]
Gut, idealerweise gibt es Aufgaben aus allen Bereichen. Es soll für jeden etwas dabei sein --- sowohl Aufgaben, bei denen man direkt die Antwort hinschreiben kann, als auch solche, bei denen man erst einmal gründlich überlegen muss.

Das ist wichtig, damit wir auch zwischen einer $5.5$ und einer $6$ differenzieren können. Es soll ja nicht so sein, dass jeder, der den Stoff einigermaßen verstanden hat, sofort die volle Punktzahl bekommt. Es muss also Aufgaben geben, bei denen man wirklich knobeln muss. Aber es gibt eben auch einfachere Aufgaben, bei denen Sie direkt etwas hinschreiben können --- wobei das natürlich auch stark davon abhängt, wie gut Sie den Stoff verstanden haben.
\end{spoken-clean}

\begin{student-interaction}[Studentenfrage]
Also allgemein, dass die Prüfung so ist, fürs schöne Aufschreiben...
\end{student-interaction}

\begin{spoken-clean}[01:28:29 - 01:29:21]
Ah, für das schöne Aufschreiben! Ja, das ist hier besonders wichtig. Das saubere Aufschreiben der Aufgaben ist ein zentraler Punkt, weil es eben darum geht zu lernen, wie man Mathematik präzise formuliert. Mehr als in anderen Vorlesungen ist eine schöne, strukturierte Lösung hier ein wesentlicher Bestandteil der Bewertung. Wie viel Zeit Sie dafür benötigen, muss ich Ihnen überlassen --- das Arbeitsverhalten ist da sehr individuell.

Gut, das wäre das Ende der Vorlesung. Ganz herzlichen Dank fürs Kommen! Es war mir eine Freude mit Ihnen, und hoffentlich sieht man sich in den folgenden Semestern Ihres Studiums mal wieder. Herzlichen Dank, einen schönen Sommer und viel Erfolg bei der Prüfung! Danke vielmals.
\end{spoken-clean}

\begin{nice-box}[Zusammenfassung der Schlüsselkonzepte]
Hier ist eine Übersicht der wichtigsten Konzepte aus der heutigen Vorlesung:
\begin{itemize}
    \item \textbf{Polynomring $K[x]$:} Die Menge der formalen Summen über einem Körper $K$. Der Grad eines Produkts verhält sich additiv ($\deg(f \cdot g) = \deg(f) + \deg(g)$), da $K$ nullteilerfrei ist.
    \item \textbf{Division mit Rest:} Für $f, g \in K[x]$ mit $f \neq 0$ existieren eindeutige $q, r$ mit $g = q \cdot f + r$ und $\deg(r) < \deg(f)$.
    \item \textbf{Irreduzibilität:} Ein nicht-konstantes Polynom ist irreduzibel, wenn es nur triviale Teiler besitzt. Irreduzible Polynome übernehmen im Polynomring die Rolle der Primzahlen in $\mathbb{Z}$.
    \item \textbf{Restklassenkörper:} Der Restklassenring $K[x]/\langle f \rangle$ ist genau dann ein Körper, wenn $f$ irreduzibel ist. Jedes Element besitzt einen eindeutigen Repräsentanten mit Grad $< \deg(f)$.
    \item \textbf{Endliche Körper:} Jeder endliche Körper hat die Ordnung $p^n$ für eine Primzahl $p$ und ist bis auf Isomorphie eindeutig bestimmt (z.\,B. $\mathbb{F}_4 \cong \mathbb{F}_2[x]/\langle x^2+x+1\rangle$).
    \item \textbf{Prüfungsstrategie:} Fokus auf präzise mathematische Argumentation, aktives Lösen von Altklausuren und Erstellen eigener Prüfungsfragen.
\end{itemize}
\end{nice-box}

% [SYSTEM] Refinement complete